\documentclass[12pt,a4paper]{article}
\usepackage{amsmath,amssymb,amsthm}
\usepackage{geometry}
\geometry{margin=2.5cm}
\newtheorem{theorem}{Theorem}[section]
\newtheorem{lemma}[theorem]{Lemma}
\newtheorem{definition}[theorem]{Definition}
\newcommand{\R}{\mathbb{R}}
\newcommand{\C}{\mathbb{C}}
\newcommand{\D}{\mathcal{D}'}
\newcommand{\Z}{\mathbb{Z}}
\newcommand{\h}{\frac{1}{2}}

\title{\bf A Geometric Criterion for the Riemann Hypothesis\\ via the Chebyshev Function}
\author{}
\date{}

\begin{document}
\maketitle

\begin{abstract}
We formulate the Riemann Hypothesis as a geometric condition on a pair of space curves
constructed from the Chebyshev function \(\psi(e^t)\).  The curves live on the unit cylinder
and interpolate a suitable half‑set of non‑trivial zeros of the Riemann zeta function.
The condition that the two curves are mirror‑symmetric forces all zeros to lie on the
critical line \(\Re s = \tfrac12\).  The construction uses only the unconditional explicit formula,
avoiding any assumption on the real parts of the zeros.  The remaining analytic task is to
prove that the mollified Chebyshev helix obtained from the prime distribution indeed
interpolates the zeta zeros in the distributional limit.
\end{abstract}

\section{Introduction}

Let \(\zeta(s)\) be the Riemann zeta function and let
\[
\rho_n = \beta_n + i\gamma_n,\qquad \gamma_n>0,\; 0<\beta_n<1,
\]
be its non‑trivial zeros, with \(\gamma_n\) in increasing order.  The functional equation
\(\zeta(s)=\zeta(1-s)\) pairs every off‑critical zero \(\beta+i\gamma\) with \(1-\beta+i\gamma\);
on the critical line \(\beta=\frac12\) a zero is self‑symmetric.

In this note we propose a geometric programme that reduces the Riemann Hypothesis (RH)
to the regularisation of a single explicit curve.  The idea is to map the critical strip
onto a cylinder of radius \(1\) and to construct two curves:
\begin{itemize}
\item \(C_+\) passes through every zero with real part \(\le \frac12\);
\item \(C_-\) is the mirror image of \(C_+\) across the plane corresponding to the
      critical line.
\end{itemize}
The curves are built directly from the prime numbers via the Chebyshev function
\(\psi(x)=\sum_{n\le x}\Lambda(n)\).  The classical explicit formula of Riemann and
von Mangoldt equates \(\psi(e^t)\) to a sum over the zeta zeros.  The rigidity of
this identity implies that the two curves \(C_\pm\) must coincide if and only if the
zero set is invariant under \(\rho\mapsto1-\rho\), i.e.\ if RH holds.

Thus the problem becomes: \emph{prove that the mollified Chebyshev helix obtained from
the prime distribution interpolates the zeta zeros in the limit where the mollifier
converges to a delta.}  We lay out the precise definitions, state the main conditional
theorem and outline the required analytic steps.

\section{Geometry of the cylinder}

We identify the critical strip \(0<\sigma<1\) with the cylinder
\(S^1\times\R\) of radius \(1\) via
\[
(\theta, t) \;=\; (2\pi\sigma,\, t).
\]
The critical line \(\sigma=\frac12\) becomes the meridian \(\theta=\pi\).
A non‑trivial zero \(\beta+i\gamma\) with \(\gamma>0\) is represented by the point
\[
P_\rho = (2\pi\beta,\; \gamma)\in S^1\times\R .
\]

Consider a smooth (or piecewise smooth) curve parametrised by height \(t\):
\[
C(t) = \bigl( \cos\Theta(t),\; \sin\Theta(t),\; t \bigr),\qquad t\in I.
\]
Its torsion, when defined, is given by the Frenet–Serret formula
\begin{equation}\label{eq:torsion}
\tau[\Theta](t) =
\frac{\Theta'\bigl(\Theta'^4 - \Theta'\Theta''' + 3\Theta''^2\bigr)}
{\bigl(1+\Theta'^2\bigr)\bigl(\Theta'^4 + \Theta''^2 + \Theta'^6\bigr)} .
\end{equation}

\section{The Chebyshev function and the explicit formula}

Let \(\psi(x)=\sum_{n\le x}\Lambda(n)\) be the Chebyshev function, with
\(\Lambda(n)=\log p\) if \(n=p^m\) and zero otherwise.
For \(t>0\) set
\[
\Psi(t) = \psi(e^t).
\]
\(\Psi(t)\) is a piecewise linear function whose distributional derivative is the
prime‑power Dirac comb
\begin{equation}\label{eq:primedist}
\Psi'(t) = \sum_{p,m} \log p\; \delta(t - m\log p)
\;=:\;\tau_0(t).
\end{equation}

The explicit formula (Riemann–von Mangoldt) relates \(\Psi\) to the zeros of \(\zeta(s)\):
\begin{equation}\label{eq:explicit}
\Psi(t) = e^t - \sum_{\rho}\frac{e^{\rho t}}{\rho}
          - \log 2\pi - \frac12\log(1-e^{-2t}) .
\end{equation}
Here the sum runs over all non‑trivial zeros \(\rho=\beta+i\gamma\) and converges in
the sense of distributions after pairing with suitable test functions.  The
archimedean terms (the last two) are smooth for \(t>0\).

\section{Construction of the two curves}

Among all non‑trivial zeros with \(\gamma>0\) we select exactly one from each
un‑ordered symmetric pair \((\beta+i\gamma, 1-\beta+i\gamma)\).  For a zero on the
critical line we keep it.  This gives a set of zeros with real parts
\(\beta_n\in(0,\frac12]\) and strictly increasing imaginary parts \(\gamma_n\).

Define a smooth approximation \(\Psi_\varepsilon = \Psi * \eta_\varepsilon\) where
\(\eta_\varepsilon(t)=\frac1\varepsilon\eta(t/\varepsilon)\) is a positive even
mollifier.  Let
\[
\Theta_\varepsilon(t) = \pi + \alpha\,\Psi_\varepsilon(t),
\]
with a normalisation constant \(\alpha\) to be fixed later.  The corresponding
smoothed curve is
\[
C_\varepsilon^+(t) = \bigl( \cos\Theta_\varepsilon(t),\; \sin\Theta_\varepsilon(t),\; t \bigr).
\]
Its mirror image is
\[
C_\varepsilon^-(t) = \bigl( \cos(2\pi-\Theta_\varepsilon(t)),\;
                         \sin(2\pi-\Theta_\varepsilon(t)),\; t \bigr).
\]

The two curves \(C_\varepsilon^\pm\) are symmetric with respect to the plane
\(\theta=\pi\).  Their torsions (with the same orientation for both) satisfy
\[
\tau[C_\varepsilon^-] = -\tau[C_\varepsilon^+].
\]
If we reverse the parametrisation of \(C_\varepsilon^-\) by setting \(s=-t\), its
torsion becomes exactly \(+\tau[C_\varepsilon^+]\) (torsion is invariant under
orientation‑preserving reparametrisation; the flip \(t\mapsto -t\) combined with
a rigid reflection yields an orientation‑preserving isometry, hence equal torsion).

\section{The geometric criterion}

The constant \(\alpha\) is chosen so that, in the limit \(\varepsilon\to0\), the
points where \(\Theta_\varepsilon(t) \equiv 0 \pmod{2\pi}\) converge to the
imaginary parts \(\gamma_n\) of the selected half‑set of zeros, and moreover
\[
\lim_{\varepsilon\to0} \Theta_\varepsilon(\gamma_n) = 2\pi\beta_n \pmod{2\pi}.
\]
That such an \(\alpha\) exists is an open analytic claim; we state it as a
conjecture.

\begin{conjecture}[Mollified interpolation]\label{conj:interp}
There exists a normalisation \(\alpha>0\) such that for the smooth helix
\(C_\varepsilon^+\) built from \(\Psi_\varepsilon\) the following holds:
\begin{enumerate}
\item \(\displaystyle\lim_{\varepsilon\to0} \tau[C_\varepsilon^+] = \tau_0\)
      in \(\D(\R)\);
\item The set \(\mathcal{Z}_\varepsilon = \{\,t : \Theta_\varepsilon(t)\equiv0
      \pmod{2\pi}\,\}\) converges as \(\varepsilon\to0\) to the imaginary parts
      \(\gamma_n\) of the half‑set of zeros, and the limiting angular values
      satisfy \(\lim_{\varepsilon\to0}\Theta_\varepsilon(t_\varepsilon)=2\pi\beta_n\).
\end{enumerate}
\end{conjecture}

The first condition is a regularisation of the torsion of the Chebyshev curve;
it is expected because \(\Psi'\) is the prime comb and the torsion formula
\eqref{eq:torsion} should tend to that comb after smoothing.  The second
condition states that the zero set of the mollified Chebyshev phase recovers
the zeta zeros.

With this conjecture, RH follows immediately.

\begin{theorem}\label{thm:RH}
Assume Conjecture~\ref{conj:interp}.  Then all non‑trivial zeros of \(\zeta(s)\)
satisfy \(\beta=\frac12\).
\end{theorem}

\begin{proof}
Consider the two limits
\[
C^+(t) = \lim_{\varepsilon\to0} C_\varepsilon^+(t),\qquad
C^-(t) = \lim_{\varepsilon\to0} C_\varepsilon^-(t)
\]
(in the sense of pointwise convergence of the zero sets and distributional
convergence of the torsion).  Both curves interpolate the same half‑set of zeros
by the second part of the conjecture.  Their torsions are identical (up to an
orientation flip that yields the same distribution \(\tau_0\)).  The explicit
formula \eqref{eq:explicit} is symmetric under replacing every zero \(\rho\) by
\(1-\rho\); therefore the prime‑built curve \(C^+\) is invariant under this
reflection.  Hence \(C^+ = C^-\), which means that for every point on the curve
\(\Theta(t) = 2\pi - \Theta(t) \pmod{2\pi}\), i.e.\ \(\Theta(t)=\pi\).
Consequently all the interpolated zeros lie on the meridian \(\theta=\pi\),
that is \(\beta_n=\frac12\) for the selected zero set.  Because the full zero set
is the union of the half‑set and its mirror, every zero satisfies \(\beta=\frac12\).
\end{proof}

\section{Towards a proof of the conjecture}

The main task is to verify the two parts of Conjecture~\ref{conj:interp}.
We indicate the required steps.

\subsection{Mollified torsion limit}

For a smooth function \(\Theta\) the torsion is given by \eqref{eq:torsion}.
Substituting \(\Theta=\pi+\alpha\Psi_\varepsilon\) and expanding in powers of
the derivatives of \(\Psi_\varepsilon\), one expects that as \(\varepsilon\to0\):
\begin{itemize}
\item \(\Psi_\varepsilon'\) converges to \(\tau_0+\text{smooth}\),
\item \(\Psi_\varepsilon''\) and \(\Psi_\varepsilon'''\) contain singular terms
      proportional to \(\delta'\) and \(\delta''\) concentrated at the prime powers.
\item The non‑linear combination in the numerator of \eqref{eq:torsion} causes
      cancellations; the dominant term is \(\Psi_\varepsilon'\) itself.
\end{itemize}
Proving that the limits of the higher derivative products vanish or combine to
leave only the Dirac comb is a problem in microlocal analysis.  The explicit form
of \(\Psi_\varepsilon\) as a convolution of a piecewise linear function gives
precise control over the singularities.

\subsection{Convergence of the zero set}

The equation \(\Theta_\varepsilon(t)=0\pmod{2\pi}\) is equivalent to
\(\Psi_\varepsilon(t) = \frac{2\pi n - \pi}{\alpha}\) for some integer \(n\).
As \(\varepsilon\to0\), the jumps of \(\Psi\) become visible: near a prime power
the mollified function spends most of its variation in an \(\varepsilon\)-neighbourhood.
The global structure of \(\Psi(t)\) is dominated by the term \(e^t\) and the sum
over zeros.  By the explicit formula \eqref{eq:explicit}, the zero crossings of
the oscillatory part \(\sum_\rho e^{\rho t}/\rho\) should correspond to the
imaginary parts \(\gamma_n\).  A careful asymptotic analysis of the convolution
integral, using the stationary phase method near the zeros, should establish the
convergence of the crossing set.

\section{Discussion}

We have presented a geometric reformulation of the Riemann Hypothesis: the prime
distribution determines a pair of curves on a cylinder whose symmetry forces the
zeta zeros onto the critical line.  The missing piece is a rigorous regularisation
argument showing that the mollified Chebyshev function indeed interpolates the
zeros in the limit.  This problem is sharply defined and lies at the interface of
analytic number theory, distribution theory, and geometry.

The programme has several attractive features:
\begin{itemize}
\item It avoids any assumption on the real parts of the zeros.
\item The construction is completely explicit, using only the Chebyshev function.
\item The symmetry that yields RH is the built‑in symmetry of the explicit formula.
\item The required limit can be investigated numerically, providing a direct
      computational test of the idea.
\end{itemize}

If the conjecture is proved, the Riemann Hypothesis follows.  Conversely, if the
mollified limit fails to interpolate the zeros unless they lie on the critical
line, we obtain a conditional proof that the existence of such a curve is
equivalent to RH.  In either case, the geometric criterion opens a new angle
on the problem.

\appendix

\section{Torsion formula}

For a curve \(\mathbf{r}(t)=(\cos\theta(t),\sin\theta(t),t)\) we compute the
Frenet–Serret torsion.  The derivatives are
\begin{align*}
\mathbf{r}' &= (-\theta'\sin\theta,\; \theta'\cos\theta,\; 1),\\
\mathbf{r}'' &= (-\theta''\sin\theta - \theta'^2\cos\theta,\;
                 \theta''\cos\theta - \theta'^2\sin\theta,\; 0),\\
\mathbf{r}''' &= (-\theta'''\sin\theta - 3\theta'\theta''\cos\theta
                  + \theta'^3\sin\theta,\;
                  \theta'''\cos\theta - 3\theta'\theta''\sin\theta
                  - \theta'^3\cos\theta,\; 0).
\end{align*}
The volume form is
\[
(\mathbf{r}'\times\mathbf{r}'')\cdot\mathbf{r}''' =
\theta' \bigl( \theta'^4 - \theta'\theta''' + 3\theta''^2 \bigr),
\]
and the squared area is
\[
|\mathbf{r}'\times\mathbf{r}''|^2 =
(1+\theta'^2)(\theta'^4 + \theta''^2 + \theta'^6).
\]
Formula \eqref{eq:torsion} follows.

\end{document}